% Generated from locale/tr/content/set-theory/story/russells-paradox-again.tex; do not hand-edit.


\olfileid[tr]{sth}{story}{rus}
\olsection{Russell Paradoksu (yeniden)}

\olref[sfr][][]{part} içinde naif bir küme teorisiyle çalıştık. Fakat
\emph{çok} naif bir anlayışa göre kümeler yalnızca yüklemlerin kaplamlarıdır.
Bu naif düşünce şu ilkeyi gerektirirdi:

\begin{defish}
  \emph{Naif Küme Oluşturma.} Her $\phi$ formülü için
  $\Setabs{x}{\phi(x)}$ vardır.
\end{defish}

Bu ilke ne kadar çekici olsa da tutarsız olduğu kanıtlanabilir. Bunu
\olref[sfr][set][rus]{sec} içinde gördük; fakat sonuç öylesine önemli ve
yalındır ki aynen yinelemeye değer.

\begin{thm}[Russell Paradoksu]
$R=\Setabs{x}{x\notin x}$ olan bir küme yoktur.
\end{thm}

\begin{proof}
$R=\Setabs{x}{x\notin x}$ varsa $R\in R$ ancak ve ancak $R\notin R$ olur;
bu bir çelişkidir.
\end{proof}

Russell bu sonucu Haziran 1901'de keşfetti. (Ancak Frege'nin
\emph{Grundgesetze}'de ana çizgileri verilen küme teorisini ele aldığından,
paradoksu tam burada sunduğumuz biçimde ifade etmedi. Buna
\olref[blv]{sec} içinde döneceğiz.) Russell 16 Haziran 1902'de Frege'ye
yazarak onun sistemindeki tutarsızlığı açıkladı. Yazışma ve arka plan için
bkz. \citet[pp.~124--8]{Heijenoort1967}.

Bu iki satırlık ispatın \emph{saf mantığın} bir sonucu olduğunu vurgulamak
gerekir. $\Setabs{x}{x\notin x}$ gösterimimizde örtük olarak (mantık dışı?)
Kaplamsallık aksiyomunu kullandığımız kabul edilebilir; çünkü
$\Setabs{x}{\phi(x)}$, varsa, $\phi$ olan nesnelerin (Kaplamsallık gereğince)
\emph{tek} kümesi demektir. Fakat sonucu şöyle ifade ederek Kaplamsallık
izleniminden bile kaçınabiliriz: \emph{elemanları tam olarak kendisinin
elemanı olmayan kümeler olan hiçbir küme yoktur}. Bunun kümelerle de pek
ilgisi yoktur. Russell'ın kendisinin gözlemlediği gibi, aynı akıl yürütme şu
sonuca götürür: \emph{Hiçbir erkek tam olarak kendini tıraş etmeyen erkekleri
tıraş etmez}. Ya da: \emph{Hiçbir pug tam olarak kendini koklamayan pugları
koklamaz}. Ve benzeri. Şematik olarak sonucun biçimi yalnızca şudur:
\[
\lnot \exists x \forall z(Rzx \liff \lnot R zz).
\]
Bu ise birinci dereceden mantığın bir teoremi (şeması)dır. Dolayısıyla küme
teorimizde ufak değişiklikler yaparak Russell Paradoksu'ndan kaçınamayız;
paradoks daha küme teorisine \emph{gelmeden önce} ortaya çıkar. (Klasik)
birinci dereceden mantığı kullanacaksak $R=\Setabs{x}{x\notin x}$ olan bir
kümenin bulunmadığını doğrudan \emph{kabul etmemiz} gerekir.

Sonuç şudur: Naif Küme Oluşturma'yı kabul edip tutarsızlıktan
\emph{kaçınmak} istiyorsanız yalnızca \emph{küme teorisini} değiştiremezsiniz.
Bunun yerine \emph{mantığınızı} baştan kurmanız gerekir.

Elbette klasik olmayan mantıklara sahip küme teorileri sunulmuştur. Fakat
bunlar en hafif deyişle standart dışıdır. Russell Paradoksu'na standart
yaklaşım, onu doğrudan bir yokluk ispatı olarak görmek ve sonra bununla nasıl
yaşanacağını öğrenmeye çalışmaktır. Biz de bu yaklaşımı izleyeceğiz.

